% Outline Project Specification Report

%The document is a report
\documentclass[12pt,a4paper]{article}

%define horizontal rule
\newcommand{\HRule}{\rule{\linewidth}{0.5mm}}

\usepackage{fullpage}

%use the listings package
\usepackage{listings}
%use the English language
\usepackage[english]{babel}
%use graphics
\usepackage{graphicx}
%use wrap figures
\usepackage{wrapfig}
%geometry stuffs
\usepackage{lscape}
%use natbib bibliography package
%\usepackage[numbers]{natbib}
%use harvard bibliography package
%\usepackage{harvard}	
%use captions
\usepackage{caption}
%use multi-row tables
\usepackage{multirow}
\usepackage{url}

\begin{document}

%use harvard citations
%\citationstyle{agsm}

%include the title page
\begin{titlepage}
 
\begin{center}

% Upper part of the page
\includegraphics[width=0.20\textwidth]{../cover_logo.png}\\[1cm]


\textsc{\LARGE Aberystwyth University}\\[0.5cm]
\textsc{\LARGE Final Report}\\[0.5cm]



 
% Title
\HRule \\[0.4cm]
{ \huge \bfseries Partridge: An Intelligent Literature Analysis and
Recommendation Suite.}\\[0.4cm]

\HRule \\[1.5cm]

 % Author and student ID
\begin{minipage}{0.4\textwidth}
\begin{flushleft} \large
\emph{Author:}\\
\textsc{James Ravenscroft}\\
jrr9@aber.ac.uk\\
090407039\\
\end{flushleft}
\end{minipage}
\begin{minipage}{0.4\textwidth}
\begin{flushright} \large
\emph{Supervisors} \\
Amanda Clare (afc)\\
Maria Liakata (mal)

\end{flushright}
\end{minipage}


\vfill

\textsc{Submitted in partial fulfilment of requirements for a Batchelor of
Science Degree at Aberystwyth University}



\vfill
 
% Bottom of the page
\textsc{\large Word Count: 0}\\
\textsc{\large Status: Draft}\\
\textsc{\large Revision: \Revision{} }\\
{\large \today}
 
\end{center}

\frontmatter
 
\end{titlepage}



%some definitions for paragraph layout stuff
\setlength{\parindent}{0pt}
\setlength{\parskip}{1.5ex plus 0.5ex minus 0.2ex}

%\tableofcontents

\pagebreak

\section{Project Description}

In academic research, there are a huge abundance of scientific papers available for
consumption for any given scientific pursuit. There are hundreds of journals,
each with thousands of articles published on a yearly basis. In 2010, the famous journal,
PLOS One (\url{http://plosone.org}), reached the 10,000 papers published
milestone \cite{PlosOne2010}. With the advent of internet publishing,
researchers have access to an immense volume of information at their disposal.

For researchers doing background reading for a paper they are authoring, or looking for 
research on a specific topic that interests them, the vast quantity of data available can often be
quite daunting and choosing which papers to read very difficult. This problem
is further exacerbated when there is a time limit on the research to be carried
out. Partridge aims to aid the reader in finding documents relative to their
interests through the use of Artificial Intelligence techniques.

The project will use a web interface to allow a user to indicate which
specific scientific fields they are interested in - either through manual entry
of interests and types of report or by allowing the application to collect a
list of papers that the user has already read and to learn their user profile.
The program will then suggest sets of papers from the corpus depending on their
contents and relevance to the user.

\section{Work to be tackled}

The main body of work to be tackled during the project is the Natural Language
Parsing (NLP) element of the application. 



\section{Project Deliverables}

\pagebreak


\bibliographystyle{IEEEannot}
\nocite{*}
\bibliography{report}

\end{document}
